其中 clear 占比为 $159/169=0.941$（脚本自动汇总）。该 fixed-site 基准刻意选在漂移主导区；K=2 奇偶粗粒化用于抑制 odd/even 微结构，使 clear 判定反映“分离时间尺度”而非 parity 伪峰。在首遇分布中 $f(0)=0$，因此 $\widetilde f(0)=f(0)=0$，检测从 $m\ge1$ 开始。与主文一致，候选峰对统一要求 $\tau-t_1\ge 2$（确保 $t_v(\tau)=\arg\min_{t_1<t<\tau}g_t$ 在开区间有定义）。该分支沿用主文同一 timescale 选峰口径：检测轨迹取 $g_m=\bar{\widetilde f}(m)$，首峰 + score 选 $t_2$（$R_{\mathrm{peak}}(\tau)=\frac{\min\{g_{t_1},g_{\tau}\}}{\max\{g_{t_1},g_{\tau}\}}$，$R_{\mathrm{peak}}=\min/\max$（即 $R_{\mathrm{peak}}(t_2)$），$R_{\mathrm{dir}}=g_{t_2}/g_{t_1}$ 仅作方向性诊断（可大于 1，不参与 phase 阈值），$\mathrm{score}(\tau)=\frac{(\tau-t_1)\,R_{\mathrm{peak}}(\tau)}{R_{\mathrm{valley}}(\tau)+10^{-12}}$）；若 $|\Delta\mathrm{score}|\le10^{-12}$，并列按“更晚 $\tau$ $\rightarrow$ 更小谷比 $\rightarrow$ 更大峰平衡比”打破；在 parity 索引上施加 $\tau-t_1\ge\mathrm{min\_sep\_pair}$，无额外 $t_{\mathrm{end}}$ 截断，并按 $t=2m$ 回映。
